\documentclass[10pt]{beamer}
\usetheme{metropolis}
\usepackage{graphicx}
\usepackage{listings}
\usepackage{booktabs}
\usepackage{xcolor}
\usepackage{hyperref}

\definecolor{codebg}{HTML}{F5F5F5}
\definecolor{codegreen}{HTML}{2E7D32}
\definecolor{codegray}{HTML}{757575}
\definecolor{codeblue}{HTML}{1565C0}

\lstdefinestyle{rstyle}{
  language=R, backgroundcolor=\color{codebg},
  basicstyle=\ttfamily\scriptsize,
  keywordstyle=\color{codeblue}\bfseries,
  stringstyle=\color{codegreen},
  commentstyle=\color{codegray}\itshape,
  breaklines=true, frame=single,
  rulecolor=\color{codegray!30}, showstringspaces=false, tabsize=2,
  morekeywords={library,ggplot,aes,geom_point,geom_line,geom_col,geom_smooth,
    geom_histogram,geom_boxplot,scale_color_manual,scale_x_log10,
    facet_wrap,theme_minimal,theme,element_text,labs,coord_flip}
}
\lstdefinestyle{pythonstyle}{
  language=Python, backgroundcolor=\color{codebg},
  basicstyle=\ttfamily\scriptsize,
  keywordstyle=\color{codeblue}\bfseries,
  stringstyle=\color{codegreen},
  commentstyle=\color{codegray}\itshape,
  breaklines=true, frame=single,
  rulecolor=\color{codegray!30}, showstringspaces=false, tabsize=4,
  morekeywords={import,from,as,pd,plt,sns,np,fig,ax,ggplot,geom_point,
    geom_line,geom_col,geom_smooth,aes,theme_minimal,labs,facet_wrap}
}

\title{Wilkinson's Grammar of Graphics: Theory}
\subtitle{Workshop 2 --- Module 1}
\author{Prof.Asoc.\ Endri Raco}
\institute{Data Visualization for Data Scientists \\ UNYT Tirana}
\date{2025/26}

\begin{document}

\begin{frame}
  \titlepage
\end{frame}

% ================================================================
\begin{frame}{Module Roadmap}
  \begin{enumerate}
    \item The problem with chart catalogues
    \item Wilkinson's seven-layer specification
    \item The algebra of graphics: data $\times$ aesthetics $\times$ geometry
    \item Layered construction: building plots incrementally
    \item Wilkinson (1999) vs Wickham (2010): theory to implementation
  \end{enumerate}
  \begin{exampleblock}{Learning Outcome}
    You will understand the Grammar of Graphics as a \emph{compositional
    system} rather than a chart catalogue, and be able to decompose any
    chart into its constituent layers.
  \end{exampleblock}
\end{frame}

% ================================================================
\section{The Problem with Chart Catalogues}
% ================================================================

\begin{frame}{Traditional Approach: Memorise Chart Types}
  \begin{center}
    \includegraphics[width=0.92\textwidth]{figures_w02m01/traditional_vs_grammar}
  \end{center}
\end{frame}

% ================================================================
\begin{frame}{Why Catalogues Fail}
  \begin{columns}[T]
    \begin{column}{0.48\textwidth}
      \textbf{The catalogue approach:}
      \begin{itemize}
        \item Finite list of named chart types
        \item Each type has its own API/function
        \item Novel combinations require new types
        \item No systematic way to modify one aspect
      \end{itemize}
    \end{column}
    \begin{column}{0.48\textwidth}
      \textbf{The grammar approach:}
      \begin{itemize}
        \item Small set of composable components
        \item Components combine freely with ``+''
        \item Novel charts emerge naturally
        \item Change one layer without rebuilding
      \end{itemize}
    \end{column}
  \end{columns}

  \vspace{6pt}
  \begin{exampleblock}{Data Insight}
    A grammar generates an \emph{infinite} set of charts from a finite
    set of rules --- exactly as a natural language grammar generates
    infinite sentences from finite vocabulary and syntax.
  \end{exampleblock}
\end{frame}

% ================================================================
\section{The Seven Layers}
% ================================================================

\begin{frame}{Wilkinson's Seven-Layer Specification}
  \begin{center}
    \includegraphics[width=0.78\textwidth]{figures_w02m01/seven_layers}
  \end{center}
\end{frame}

% ================================================================
\begin{frame}{Layer 1: Data + Aesthetics}
  \begin{columns}[T]
    \begin{column}{0.55\textwidth}
      The foundation: map \textbf{data variables} to \textbf{visual channels}.

      \vspace{4pt}
      \centering
      \begin{tabular}{ll}
        \toprule
        \textbf{Variable} & \textbf{Channel} \\
        \midrule
        GDP          & x position \\
        Life Expectancy & y position \\
        Continent    & colour (hue) \\
        Population   & size \\
        Year         & animation frame \\
        \bottomrule
      \end{tabular}
    \end{column}
    \begin{column}{0.42\textwidth}
      \begin{alertblock}{Pro-Tip}
        The aesthetic mapping is the single
        most important decision. A chart is
        only as good as its variable-to-channel
        assignment. Revisit Cleveland \& McGill
        (W01-M03) to choose the best channel.
      \end{alertblock}
    \end{column}
  \end{columns}
\end{frame}

% ================================================================
\begin{frame}{Layer 2: Geometries}
  \textbf{Geometries} are the visual marks that represent data:

  \vspace{4pt}
  \centering
  \begin{tabular}{llp{4.5cm}}
    \toprule
    \textbf{Geom} & \textbf{Mark} & \textbf{Best For} \\
    \midrule
    \texttt{geom\_point}     & Circle/dot    & Relationships (scatter) \\
    \texttt{geom\_line}      & Connected line & Temporal trends \\
    \texttt{geom\_col/bar}   & Rectangle     & Comparisons \\
    \texttt{geom\_histogram} & Binned rect.  & Distributions \\
    \texttt{geom\_boxplot}   & Box + whisker & Group summaries \\
    \texttt{geom\_area}      & Filled region & Cumulative / stacked \\
    \texttt{geom\_tile}      & Coloured cell & Heatmaps \\
    \bottomrule
  \end{tabular}

  \vspace{6pt}
  \begin{exampleblock}{Data Insight}
    The geometry is \emph{not} the chart type --- it is one layer of the
    chart. Multiple geometries can be stacked: \texttt{geom\_point() +
    geom\_smooth()} is two layers on one axes.
  \end{exampleblock}
\end{frame}

% ================================================================
\begin{frame}{Same Data, Different Geometry}
  \begin{center}
    \includegraphics[width=0.95\textwidth]{figures_w02m01/same_data_diff_geom}
  \end{center}
  \begin{alertblock}{Pro-Tip}
    The data and aesthetic mapping are identical across all three panels.
    Only the geometry changes. This is the power of the grammar: swap one
    layer without rewriting anything else.
  \end{alertblock}
\end{frame}

% ================================================================
\begin{frame}{Layers 3--5: Scales, Statistics, Facets}
  \begin{columns}[T]
    \begin{column}{0.32\textwidth}
      \textbf{3. Scales}

      Map data values to visual
      properties:

      \begin{itemize}
        \item Continuous $\rightarrow$ position
        \item Discrete $\rightarrow$ colour hue
        \item Log transform
        \item Date formatting
      \end{itemize}
    \end{column}
    \begin{column}{0.32\textwidth}
      \textbf{4. Statistics}

      Transform data before
      rendering:

      \begin{itemize}
        \item \texttt{stat\_bin} (histogram)
        \item \texttt{stat\_smooth} (LOESS)
        \item \texttt{stat\_summary} (mean)
        \item \texttt{stat\_count} (bar)
      \end{itemize}
    \end{column}
    \begin{column}{0.32\textwidth}
      \textbf{5. Facets}

      Split data into panels
      by a categorical variable:

      \begin{itemize}
        \item \texttt{facet\_wrap}
        \item \texttt{facet\_grid}
        \item Shared or free scales
        \item Tufte's small multiples
      \end{itemize}
    \end{column}
  \end{columns}
\end{frame}

% ================================================================
\begin{frame}{Layers 6--7: Coordinates \& Themes}
  \begin{columns}[T]
    \begin{column}{0.48\textwidth}
      \textbf{6. Coordinate System}

      Defines the geometry of the plane:

      \begin{itemize}
        \item \texttt{coord\_cartesian} (default)
        \item \texttt{coord\_flip} (swap x/y)
        \item \texttt{coord\_polar} (pie, coxcomb)
        \item \texttt{coord\_sf} (map projection)
      \end{itemize}

      \vspace{4pt}
      {\small A bar chart in Cartesian is bars;
      the same bar chart in polar is a pie chart.
      Only the coordinate system changed.}
    \end{column}
    \begin{column}{0.48\textwidth}
      \textbf{7. Theme}

      Controls all non-data visual elements:

      \begin{itemize}
        \item Fonts: family, size, weight
        \item Backgrounds: panel, plot, strip
        \item Gridlines: major, minor
        \item Legend: position, box, key
      \end{itemize}

      \vspace{4pt}
      {\small Themes are separable from the data
      encoding --- change the theme without
      touching the analysis.}
    \end{column}
  \end{columns}
\end{frame}

% ================================================================
\begin{frame}{Same Data, Different Aesthetic Strategy}
  \begin{center}
    \includegraphics[width=0.92\textwidth]{figures_w02m01/same_data_diff_aes}
  \end{center}
  \begin{exampleblock}{Data Insight}
    Colour and faceting are both valid ways to encode a grouping
    variable. Colour keeps all data on one panel (good for $\leq$3
    groups); faceting separates panels (better for $>$3 groups or
    when groups overlap heavily).
  \end{exampleblock}
\end{frame}

% ================================================================
\section{The Algebra of Graphics}
% ================================================================

\begin{frame}{Data $\times$ Aesthetics $\times$ Geometry = Chart}
  \begin{center}
    \includegraphics[width=0.92\textwidth]{figures_w02m01/algebra}
  \end{center}
\end{frame}

% ================================================================
\begin{frame}{Wilkinson's Specification Pipeline}
  \begin{center}
    \includegraphics[width=0.95\textwidth]{figures_w02m01/specification_pipeline}
  \end{center}
  \begin{exampleblock}{Data Insight}
    Wilkinson's pipeline is a formal specification: Variables $\rightarrow$
    Algebra $\rightarrow$ Scales $\rightarrow$ Statistics $\rightarrow$
    Geometry $\rightarrow$ Coordinates $\rightarrow$ Aesthetics (rendering).
    Each step is independent and can be modified without affecting the others.
  \end{exampleblock}
\end{frame}

% ================================================================
\section{Layered Construction}
% ================================================================

\begin{frame}{Building a Plot Layer by Layer}
  \begin{center}
    \includegraphics[width=0.95\textwidth]{figures_w02m01/layering}
  \end{center}
  \begin{alertblock}{Pro-Tip}
    Start with data + axes (see the variable ranges), add geometry
    (see the data), add statistics (see the pattern), then polish with
    labels and theme. This ``progressive reveal'' mirrors the design
    process from W01-M06.
  \end{alertblock}
\end{frame}

% ================================================================
\begin{frame}[fragile]{Layered Construction in R (ggplot2)}
  \begin{columns}[T]
    \begin{column}{0.52\textwidth}
\begin{lstlisting}[style=rstyle]
library(ggplot2)

# Layer 1: data + aesthetics
ggplot(mpg,
  aes(x = displ, y = hwy,
      color = class)) +

# Layer 2: geometry
  geom_point(size = 2,
             alpha = 0.6) +

# Layer 3: statistic
  geom_smooth(method = "lm",
              se = TRUE,
              color = "grey40") +

# Layer 4: scale
  scale_color_brewer(
    palette = "Set2") +

# Layer 5: labels + theme
  labs(title = "Displacement vs MPG",
       x = "Engine (L)",
       y = "Highway MPG") +
  theme_minimal(base_size = 11)
\end{lstlisting}
    \end{column}
    \begin{column}{0.45\textwidth}
      \includegraphics[width=\textwidth]{figures_w02m01/r_layered}

      \vspace{4pt}
      {\small Each \texttt{+} adds one layer.
      The order of layers determines
      the drawing order (first layer
      drawn first, behind later layers).}
    \end{column}
  \end{columns}
\end{frame}

% ================================================================
\begin{frame}[fragile]{Layered Construction in Python (plotnine)}
  \begin{columns}[T]
    \begin{column}{0.52\textwidth}
\begin{lstlisting}[style=pythonstyle]
from plotnine import *
import pandas as pd

# plotnine uses identical syntax
# to R's ggplot2

(
  ggplot(mpg,
    aes(x="displ", y="hwy",
        color="class"))

  # Layer 2: geometry
  + geom_point(size=2, alpha=0.6)

  # Layer 3: statistic
  + geom_smooth(method="lm",
                se=True,
                color="grey")

  # Layer 4: scale
  + scale_color_brewer(
      type="qual",
      palette="Set2")

  # Layer 5: labels + theme
  + labs(title="Displacement vs MPG",
         x="Engine (L)",
         y="Highway MPG")
  + theme_minimal()
)
\end{lstlisting}
    \end{column}
    \begin{column}{0.45\textwidth}
      \includegraphics[width=\textwidth]{figures_w02m01/py_plotnine}

      \vspace{4pt}
      {\small \texttt{plotnine} is a near-exact
      port of ggplot2 to Python. The
      only syntax difference: strings
      for column names instead of bare
      symbols, and \texttt{+} at the start
      of each line (Python line
      continuation).}
    \end{column}
  \end{columns}
\end{frame}

% ================================================================
\section{Wilkinson vs Wickham}
% ================================================================

\begin{frame}{Theory (1999) vs Implementation (2010)}
  \begin{center}
    \includegraphics[width=0.85\textwidth]{figures_w02m01/wilkinson_vs_wickham}
  \end{center}
\end{frame}

% ================================================================
\begin{frame}{Wickham's Key Contributions}
  Wickham's 2010 paper and the \texttt{ggplot2} package extended
  Wilkinson's theory with four practical innovations:

  \begin{enumerate}
    \item \textbf{Layered architecture}: multiple geom/stat pairs stacked
          on a single axes via \texttt{+}, each with its own aesthetic
          mapping if needed
    \item \textbf{Position adjustments}: \texttt{dodge}, \texttt{stack},
          \texttt{fill}, \texttt{jitter} --- modify the placement of
          marks without changing the geometry
    \item \textbf{Default statistics}: each geom has a default stat
          (\texttt{geom\_bar} defaults to \texttt{stat\_count}), reducing
          boilerplate
    \item \textbf{Theme system}: a complete separation of data encoding
          from visual styling, enabling one-line theme swaps
  \end{enumerate}

  \begin{alertblock}{Pro-Tip}
    Read Wickham (2010) as the ``user manual'' for the grammar. Read
    Wilkinson (1999/2005) for the ``why'' behind the design.
  \end{alertblock}
\end{frame}

% ================================================================
\section{Synthesis}
% ================================================================

\begin{frame}{Decomposing Any Chart into Grammar Components}
  \textbf{Exercise}: decompose a grouped bar chart into its layers.

  \vspace{4pt}
  \centering
  \begin{tabular}{ll}
    \toprule
    \textbf{Layer} & \textbf{Specification} \\
    \midrule
    Data        & sales data frame (product, quarter, revenue) \\
    Aesthetics  & x = quarter, y = revenue, fill = product \\
    Geometry    & \texttt{geom\_col(position = "dodge")} \\
    Scale       & \texttt{scale\_fill\_brewer(palette = "Set2")} \\
    Statistics  & identity (raw values, no transformation) \\
    Facets      & none (single panel) \\
    Coordinates & Cartesian (default) \\
    Theme       & \texttt{theme\_minimal()} \\
    \bottomrule
  \end{tabular}

  \vspace{6pt}
  {\small Changing any single row produces a different chart. Change
  geometry to \texttt{geom\_line()} $\rightarrow$ line chart. Change
  coordinates to \texttt{coord\_polar()} $\rightarrow$ radial bar.
  Change facets to \texttt{facet\_wrap(\textasciitilde{}product)}
  $\rightarrow$ small multiples.}
\end{frame}

% ================================================================
\begin{frame}{Key Takeaways}
  \begin{enumerate}
    \item The Grammar of Graphics replaces a \textbf{chart catalogue}
          with a \textbf{compositional system}: data + aesthetics +
          geometry + scales + statistics + facets + coordinates + theme.
    \item The \textbf{algebra} Data $\times$ Aesthetics $\times$ Geometry
          generates infinite chart types from finite components.
    \item \textbf{Layered construction} builds plots incrementally:
          each \texttt{+} adds one component. Start with data, end with theme.
    \item The same data and aesthetics with a \textbf{different geometry}
          produce a different chart --- swap one layer, keep the rest.
    \item \textbf{ggplot2} (R) and \textbf{plotnine} (Python) implement
          the grammar with near-identical syntax.
  \end{enumerate}
\end{frame}

% ================================================================
\begin{frame}{Essential Readings for This Module}
  \begin{itemize}
    \item Wilkinson, L. (2005). \emph{The Grammar of Graphics}, 2nd ed.,
          Chapters 1--2 (Introduction, How to Make a Pie).
    \item Wickham, H. (2010). ``A Layered Grammar of Graphics.''
          \emph{Journal of Computational and Graphical Statistics}, 19(1), 3--28.
    \item Wickham, H. (2016). \emph{ggplot2: Elegant Graphics for Data
          Analysis}, 2nd ed., Chapter 1 (Getting Started).
    \item Wilke, C.~O. (2019). \emph{Fundamentals of Data Visualization},
          Chapter 2 (Visualizing Data: Mapping Data onto Aesthetics).
  \end{itemize}

  \vspace{6pt}
  \textbf{Next Module}: Wickham's Layered Grammar \& ggplot2 (W02-M02)
\end{frame}

\end{document}
